\documentclass[11pt]{article}
\usepackage{fontspec}
\usepackage[french,english]{babel}
\usepackage{amsmath}
\usepackage{amssymb}
\usepackage{amsthm}
\usepackage{mathrsfs}
\usepackage{tikz-cd}
\usepackage{enumitem}

\newcommand{\printedpage}[1]{\marginpar{\footnotesize p.~#1}}
\newcommand{\msp}[1]{\marginpar{\footnotesize [#1]}}
\newcommand{\Ho}{\mathrm{Ho}}
\newcommand{\lc}{\mathrm{lc}}
\newcommand{\uFib}{\underline{\mathrm{Fib}}}

\begin{document}

\setcounter{footnote}{131}

\printedpage{160}
% head: A. Grothendieck : Dérivateurs, Ch. XII, Édition des Universités de Montpellier II et Paris VII
% foot: Version 30 mars 2003

\msp{page 247}

\textbf{Corollary.} {\itshape Description of the functor
\[
(\xi_1,\xi_2)\longmapsto\xi_1\times_S\xi_2 : \Ho^{\lc}_W(S)\times\Ho^{\lc}_W(S)\longrightarrow\Ho^{\lc}_W(S).
\]
This functor really is a cartesian product in} $\Ho^{\lc}_W(S)$ (\footnote{but this is a special case of the cases of $\Ho_W$ (if $\mathcal{M}$ has binary products), replacing $\mathcal{M}$ by $\mathcal{M}/S$.}), {\itshape where finite products exist} (\footnote{Say also that the functors $f^* : \Ho^{\lc}_W(S)\longrightarrow\Ho^{\lc}_W(S')$ (for $f : S'\longrightarrow S$) commute with finite products.}).

To prove this last point, setting
\[
\mathcal{H}=\Ho^{\lc}_W(S),
\]
we must show that the functors
\[
\begin{tikzcd}[ampersand replacement=\&]
\mathcal{H} \arrow[r, shift left, "\delta"] \& \mathcal{H}\times\mathcal{H} \arrow[l, shift left, "\pi"]
\end{tikzcd}
\]
are adjoint, with
\[
\begin{array}{rclcrcl}
\delta & = & \mathrm{diag}_{\mathcal{H}} & : & \xi & \longmapsto & (\xi,\xi)\\
\pi & = & \text{``product''} & : & (\xi_1,\xi_2) & \longmapsto & \xi_1\times_S\xi_2.
\end{array}
\]
We are going to define the adjunction homomorphisms
\[
\begin{array}{rcllrcl}
\mathrm{id}_{\mathcal{H}} & \xrightarrow{\alpha} & \pi\delta & \text{i.e.} & \xi & \xrightarrow{\alpha_\xi} & \xi\times_S\xi\\
\mathrm{id}_{\mathcal{H}\times\mathcal{H}} & \xleftarrow{\beta} & \delta\pi & \text{i.e.} & (\xi_1,\xi_2) & \xleftarrow{\beta_{\xi_1,\xi_2}} & (\xi_1\times_S\xi_2,\xi_1\times_S\xi_2).
\end{array}
\]
For this, we describe $\mathcal{H}$ as
\[
\mathcal{H}=(\uFib_W S)(W_S)^{-1}
\]
and we note that $\delta$, $\pi$ come from adjoint functors at the level of $\mathcal{H}_0=\uFib_W S$
\[
\begin{tikzcd}[ampersand replacement=\&]
\mathcal{H}_0 \arrow[r, shift left, "\delta_0"] \& \mathcal{H}_0\times\mathcal{H}_0 \arrow[l, shift left, "\pi_0"]
\end{tikzcd},
\]
which give rise to $\alpha_0$, $\beta_0$
\[
\begin{array}{rcl}
\mathrm{id}_{\mathcal{H}_0} & \xrightarrow{\alpha_0} & \pi_0\delta_0\\
\mathrm{id}_{\mathcal{H}_0\times\mathcal{H}_0} & \xleftarrow{\beta_0} & \delta_0\pi_0.
\end{array}
\]

\printedpage{161}
% head: A. Grothendieck : Dérivateurs, Ch. XII, Édition des Universités de Montpellier II et Paris VII
% foot: Version 30 mars 2003

\msp{page 248}

We define $\alpha$, $\beta$ by passage to the categories of fractions for $\alpha_0$, $\beta_0$. The 2 compatibilities, being satisfied for $\alpha_0$, $\beta_0$, are also satisfied for $\alpha$, $\beta$. Q.e.d.

I single out here the

\textbf{Lemma.} {\itshape Let \begin{tikzcd}[ampersand replacement=\&]
\mathcal{X} \arrow[r, shift left, "f"] \& \mathcal{Y} \arrow[l, shift left, "g"]
\end{tikzcd} be a pair of adjoint functors, and let $\Sigma$, $\Sigma'$ be localizers in $\mathcal{X}$, $\mathcal{Y}$ such that $f(\Sigma)\subseteq\Sigma'$, $g(\Sigma')\subseteq\Sigma$, so that we have a pair of functors
\[
\begin{tikzcd}[ampersand replacement=\&]
\mathcal{X}\Sigma^{-1} \arrow[r, shift left, "\bar{f}"] \& \mathcal{Y}\Sigma'^{-1} \arrow[l, shift left, "\bar{g}"]
\end{tikzcd}.
\]
Then $\bar{f}$, $\bar{g}$ are adjoint, the adjunction arrows being deduced from those for $f$, $g$.}

\textbf{Corollary.} {\itshape Let $(\mathcal{M},W)$ be a category of models, suppose that in $\mathcal{M}$ binary products (resp. finite products) exist, and that for $u_1 : X_1\longrightarrow X'_1$ and $u_2 : X_2\longrightarrow X'_2$, if $u_1,u_2\in W$, then $u_1\times u_2 : X_1\times X_2\longrightarrow X'_1\times X'_2$ is in $W$. Then in $\mathcal{M}W^{-1}$ binary products (resp. finite products) exist, and the localization functor $\mathcal{M}\longrightarrow\mathcal{M}W^{-1}$ commutes with them.}

We apply the lemma to
\[
\begin{tikzcd}[ampersand replacement=\&]
\mathcal{M} \arrow[r, shift left, "\delta"] \& \mathcal{M}\times\mathcal{M} \arrow[l, shift left, "\pi"]
\end{tikzcd}
\qquad \text{and } W,\ W\times W,
\]
and for the final object,
\[
\begin{tikzcd}[ampersand replacement=\&]
\mathcal{M} \arrow[r, shift left, "e"] \& e \arrow[l, shift left, "\bullet"]
\end{tikzcd}
\qquad\qquad \text{and } W,\ W_0=\{\mathrm{id}_e\}.
\]

\msp{page 249}

I am now going to construct the functor
\[
i_* : \Ho_W(\Psi)\longrightarrow\Ho_W(\tilde{\Psi}).
\]
I must first prove the

\textbf{Lemma 1.} {\itshape Let $\mathcal{F}$ be the full subcategory of $\mathcal{M}(\Psi)$, category of diagrams \begin{tikzcd}[ampersand replacement=\&, row sep=small, column sep=small]
 \& X_1 \arrow[dl, "p_1"'] \\
S \& \\
 \& X_2 \arrow[ul, "p_2"]
\end{tikzcd} in $S$ \texttt{[rather in }$\mathcal{M}$\texttt{, over }$S$\texttt{]}, consisting of those for which $p_1$, $p_2$ are in Fib$_W$. Let $W_{\mathcal{F}}$ be the localizer induced on $\mathcal{F}$ by $W(\Psi)$ (termwise $W$-equivalences). Then
\[
\mathcal{F}W_{\mathcal{F}}^{-1}\xrightarrow{\sim}\mathcal{M}(\Psi)W(\Psi)^{-1}\ (\overset{\text{def}}{=}\Ho_W(\Psi)).
\]}

\printedpage{162}
% head: A. Grothendieck : Dérivateurs, Ch. XII, Édition des Universités de Montpellier II et Paris VII
% foot: Version 30 mars 2003

We are going to prove this lemma below. Assuming it, to construct $i_*$, it suffices to construct
\[
\mathcal{F}\longrightarrow\Ho_W(\tilde{\Psi})
\]
which maps $W_{\mathcal{F}}$ to isomorphisms. I shall take, of course,
\[
\begin{gathered}
\begin{tikzcd}[ampersand replacement=\&, row sep=small, column sep=small]
 \& X_1 \arrow[dl] \& \& \& \& X_1 \arrow[dl] \& \\
S \& \& \longmapsto \& \& S \& \& X_1\times_S X_2 \arrow[ul] \arrow[dl] \\
 \& X_2 \arrow[ul] \& \& \& \& X_2 \arrow[ul] \&
\end{tikzcd}\\
\mathcal{F}\quad\longrightarrow\quad\Ho_W(\Psi).
\end{gathered}
\] % [sic] \Ho_W(\Psi) -> \Ho_W(\tilde{\Psi}), the image is a diagram of type \tilde{\Psi} and the functor to be constructed is \mathcal{F}\longrightarrow\Ho_W(\tilde{\Psi})
Let us check that it transforms arrows of $W_{\mathcal{F}}$ into isomorphisms.

\msp{page 250}

\[
\begin{tikzcd}[ampersand replacement=\&, row sep=small, column sep=small]
 \& X_1\times_S X_2 \arrow[dl] \arrow[dd] \& \& X'_1\times_{S'}X'_2 \arrow[ll, "f_3"'] \arrow[dl] \arrow[dd] \\
X_1 \arrow[dd] \& \& X'_1 \arrow[ll, crossing over, "f_1"'] \& \\
 \& X_2 \arrow[dl] \& \& X'_2 \arrow[ll, "f_2"'] \arrow[dl] \\
S \& \& S' \arrow[ll, "f_0"'] \arrow[from=uu, crossing over] \&
\end{tikzcd}
\]

We must prove that if $f_0,f_1,f_2\in W$, then $f_3\in W$. But let us consider the inverse image of the cartesian square $Q$ under $f_0 : S'\longrightarrow S$, whence a cartesian square $Q_{S'}$ and a factorization
\[
\begin{tikzcd}[ampersand replacement=\&]
Q \& Q_{S'} \arrow[l, "\tilde{f}"'] \& Q' \arrow[l, "g"'] \arrow[ll, bend left, "f"]
\end{tikzcd}
\]
of $Q'\longrightarrow Q$. Since $f_0\in W$, and since $X_1$, $X_2$, $X_1\times_S X_2$ are in $\uFib_W S$, the horizontal arrows realizing $Q_{S'}\longrightarrow Q$ are all in $W$. This, plus the fact that $f_1,f_2\in W$, imply that in the morphism $g$, the arrows $g_1$, $g_2$ are in $W$. It follows that the arrow
\[
g_3 : X'_1\times_{S'}X'_2\longrightarrow(X_1\times_S X_2)\times_S S'
\]
is as well (theory of the product in $\uFib_W S'$ and $\Ho^{\lc}_W(S)$), hence $f_3=\tilde{f}_3g_3\in W$, q.e.d.

\msp{page 251}

Let us note that we have a canonical isomorphism
\[
\alpha : i^*i_*\xrightarrow{\sim}\mathrm{id}_{\Ho_W(\Phi)}. \tag*{$(*)$}
\]
I feel like proving that the functors
\[
\begin{tikzcd}[ampersand replacement=\&]
\Ho_W(\Phi) \arrow[r, shift right, "i_*"'] \& \Ho_W(\tilde{\Phi}) \arrow[l, shift right, "i^*"']
\end{tikzcd}
\]
\printedpage{163}
% head: A. Grothendieck : Dérivateurs, Ch. XII, Édition des Universités de Montpellier II et Paris VII
% foot: Version 30 mars 2003
are adjoint, and that $(*)$ is one of the adjunction arrows. We must define the other adjunction arrow
\[
\beta : i_*i^*\longleftarrow\mathrm{id}_{\Ho_W(\tilde{\Phi})}. \tag*{$(**)$}
\]
But for this, I must use the

\textbf{Lemma 2.} Let $\mathcal{G}$ be the full subcategory of $\mathcal{M}(\tilde{\Phi})$, consisting of the \begin{tikzcd}[ampersand replacement=\&]
X_1 \arrow[d, "p_1"'] \& X_3 \arrow[l, "q_2"'] \arrow[d, "q_1"] \\
S \& X_2 \arrow[l, "p_2"]
\end{tikzcd} such that $p_1,p_2\in\mathrm{Fib}_W$ and let $W_{\mathcal{G}}$ be the localizer induced by $W(\tilde{\Phi})$ (termwise $W$-equivalences). Then the inclusion $\mathcal{G}\longrightarrow\mathcal{M}(\tilde{\Phi})$ induces an equivalence
\[
\mathcal{G}W_{\mathcal{G}}^{-1}\xrightarrow{\sim}\mathcal{M}(\tilde{\Phi})W(\tilde{\Phi})^{-1}\overset{\text{def}}{=}\Ho_W(\tilde{\Phi}).
\]

I shall for the moment assume this

\msp{page 252}

lemma, just like lemma 1 (p. 249). I am going to consider $i^*$, $i_*$ as arrows between the categories $\mathcal{F}W_{\mathcal{F}}^{-1}$, $\mathcal{G}W_{\mathcal{G}}^{-1}$. They are induced by the arrows % [sic] categories -> catégories, missing accent
\[
\begin{tikzcd}[ampersand replacement=\&]
\mathcal{F} \arrow[r, shift right, "i_{0\,*}"'] \& \mathcal{G} \arrow[l, shift right, "i_0^*"']
\end{tikzcd}
\]
between $\mathcal{F}$ and $\mathcal{G}$ themselves, % [sic] eux mêmes -> eux-mêmes, missing hyphen
\[
\begin{tikzcd}[ampersand replacement=\&]
 \& X_1 \arrow[dd] \& X_3 \arrow[l] \arrow[dd] \& \& X_1 \arrow[dd] \& \\
i_0^* :\quad \mathcal{G}\quad\longrightarrow\quad\mathcal{F} : \& \& \& \longmapsto \& \& \\
 \& S \& X_2 \arrow[l] \& \& S \& X_2 \arrow[l]
\end{tikzcd}
\]
\[
\begin{tikzcd}[ampersand replacement=\&]
 \& X_1 \arrow[dd] \& \& \& X_1 \arrow[dd] \& X_1\times_S X_2 \arrow[l] \arrow[dd] \\
i_{0\,*} :\quad \mathcal{F}\quad\longrightarrow\quad\mathcal{G} : \& \& \& \longmapsto \& \& \\
 \& S \& X_2 \arrow[l] \& \& S \& X_2. \arrow[l]
\end{tikzcd}
\]
These two functors are adjoint to one another, and compatible with the localizers $W_{\mathcal{F}}$, $W_{\mathcal{G}}$. Thus the functors deduced from $i_0^*$, $i_{0\,*}$ by passage to the categories of fractions are also adjoint (cf. lemma p. 248).

Thus, modulo the proof of lemmas 1 and 2, we have proved this:

\printedpage{164}
% head: A. Grothendieck : Dérivateurs, Ch. XII, Édition des Universités de Montpellier II et Paris VII
% foot: Version 30 mars 2003

\msp{page 253}

\textbf{Theorem.} {\itshape Let $(\mathcal{M},W)$ be a category of models ($W\subseteq\mathit{Fl}\,\mathcal{M}$) satisfying the left factorization axiom (p. 240), thus giving rise to a good theory of $W$-homotopic fibers. Let
\[
\begin{tikzcd}[ampersand replacement=\&, row sep=small, column sep=small]
 \& 1 \& \& \& \& 1 \arrow[dr] \& \\
\Psi=0 \arrow[ur] \arrow[dr] \& \& {} \arrow[r, "i"] \& {} \& \tilde{\Psi}=0 \arrow[ur] \arrow[dr] \& \& 3 \\
 \& 2 \& \& \& \& 2 \arrow[ur] \&
\end{tikzcd}
\]
(open immersions in} $\mathrm{Cat}${\itshape ), and let us consider the restriction functor
\[
i^* : \Ho_W(\tilde{\Psi})\longrightarrow\Ho_W(\Psi).
\]
This functor admits a right adjoint
\[
i_* : \Ho_W(\Psi)\longrightarrow\Ho_W(\tilde{\Psi}).
\]
This functor $i_*$ is ``computed'' as follows: its value at \begin{tikzcd}[ampersand replacement=\&]
X_1 \arrow[d, "p_1"'] \& \\
S \& X_2 \arrow[l, "p_2"]
\end{tikzcd} is canonically isomorphic (in} $\Ho_W(\tilde{\Psi})${\itshape ) to \begin{tikzcd}[ampersand replacement=\&]
X_1 \arrow[d] \& X_1\times_S X_2 \arrow[l] \arrow[d] \\
S \& X_2 \arrow[l]
\end{tikzcd}, if we suppose that $p_1$ or $p_2$ is a $W$-fibration. In the general case it is computed by factoring $p_1$ as
\[
X_1\xrightarrow{i_1}\bar{X}_1\xrightarrow{p'_1}S,\qquad i_1\in W,\ p'_1\in\mathit{Fib}_W,
\]
and by taking \begin{tikzcd}[ampersand replacement=\&]
\bar{X}_1 \arrow[d] \& \bar{X}_1\times_S X_2 \arrow[l] \arrow[d] \\
S \& X_2 \arrow[l]
\end{tikzcd}, or by

\msp{page 254}

proceeding symmetrically, or by simultaneously factoring $p_1$ and $p_2$, and taking % [sic] simultanement -> simultanément, missing accent
\[
\begin{tikzcd}[ampersand replacement=\&]
\bar{X}_1 \arrow[d] \& \bar{X}_1\times_S\bar{X}_2 \arrow[l] \arrow[d] \\
S \& \bar{X}_2. \arrow[l]
\end{tikzcd}
\]}

\end{document}
